% ============================================================================
% Chapter 5 — Components
% ============================================================================
\chapter{Components}

This chapter lists all built-in component types. Components are plain structs
with no virtual functions. They hold data only; logic lives in systems
(free functions that call \texttt{world.forEach}).

\textbf{Headers:}
\begin{itemize}[noitemsep]
  \item \texttt{ecs/Components.hpp} — 2D components
  \item \texttt{ecs/Components3D.hpp} — 3D transform components
  \item \texttt{ecs/CameraComponents.hpp} — camera settings
  \item \texttt{ecs/LightComponents.hpp} — lighting
  \item \texttt{ecs/MeshComponents.hpp} — mesh rendering
\end{itemize}

All components are in namespace \texttt{Caffeine::ECS}.

% ----------------------------------------------------------------------------
\section{Transform}

Spatial state in 3D space. Used for all movable objects, both 2D and 3D.

\begin{lstlisting}
struct Transform {
    Vec3 position = {0, 0, 0};
    Vec3 rotation = {0, 0, 0};  // Euler angles in radians
    Vec3 scale    = {1, 1, 1};
};
\end{lstlisting}

For 2D games, use \texttt{position.x}, \texttt{position.y}, and
\texttt{rotation.z} (yaw). Leave \texttt{position.z} at 0.

\begin{lstlisting}
auto& t = world.add<Transform>(e);
t.position = {320.0f, 240.0f, 0.0f};
t.rotation.z = 1.5707f;   // 90 degrees
t.scale = {2.0f, 2.0f, 1.0f};
\end{lstlisting}

% ----------------------------------------------------------------------------
\section{Velocity2D and Acceleration2D}

\begin{lstlisting}
struct Velocity2D {
    f32 x = 0.0f;
    f32 y = 0.0f;
};

struct Acceleration2D {
    f32 x = 0.0f;
    f32 y = 0.0f;
};
\end{lstlisting}

These are plain data containers. Your movement system reads them and writes
to \texttt{Transform}. Alternatively, use the built-in Physics2D system
(Chapter~\ref{chap:physics2d}), which manages velocity integration
automatically when a \texttt{RigidBody2D} is present.

% ----------------------------------------------------------------------------
\section{Sprite}

Associates an entity with a named sprite and a frame index for sprite-sheet
animation.

\begin{lstlisting}
struct Sprite {
    std::string name;        // asset name in the asset manager
    u32         frameIndex = 0;
};
\end{lstlisting}

\begin{lstlisting}
world.add<Sprite>(e, Sprite{ "player_idle", 0 });
\end{lstlisting}

% ----------------------------------------------------------------------------
\section{Health}

\begin{lstlisting}
struct Health {
    u32 current = 100;
    u32 max     = 100;
};
\end{lstlisting}

\begin{lstlisting}
Health& hp = world.add<Health>(e, Health{100, 100});

// take damage
hp.current = (damage >= hp.current) ? 0 : hp.current - damage;

if (hp.current == 0) {
    world.destroy(e);
}
\end{lstlisting}

% ----------------------------------------------------------------------------
\section{Tag and DisabledTag}

\begin{lstlisting}
struct Tag {};         // custom marker — add a using alias
struct DisabledTag {}; // entity is excluded from most systems
\end{lstlisting}

Tags are zero-size structs. Use them as type-safe flags:

\begin{lstlisting}
// Define project-specific tags
using PlayerTag  = Caffeine::ECS::Tag;
using EnemyTag   = Caffeine::ECS::Tag;   // use separate structs per project

// Disable an entity temporarily
world.add<DisabledTag>(e);

// Re-enable
world.remove<DisabledTag>(e);
\end{lstlisting}

% ----------------------------------------------------------------------------
\section{PersistentComponent}

Marks an entity to survive scene transitions (equivalent to
\texttt{DontDestroyOnLoad} in other engines).

\begin{lstlisting}
struct PersistentComponent {
    bool dontDestroyOnLoad = false;
};
\end{lstlisting}

\begin{lstlisting}
world.add<PersistentComponent>(musicEntity, PersistentComponent{true});
\end{lstlisting}

% ----------------------------------------------------------------------------
\section{ParticleEmitterComponent}

Drives a CPU particle system attached to an entity.

\begin{lstlisting}
struct ParticleEmitterComponent {
    int   maxParticles = 100;
    f32   emissionRate = 10.0f;   // particles per second
    f32   lifetime     = 1.0f;    // seconds per particle
    Vec2  velocityMin  = {-1.0f, -1.0f};
    Vec2  velocityMax  = {1.0f,  1.0f};
    u32   startColor   = 0xFFFFFFFF;
    u32   endColor     = 0x00FFFFFF;   // RGBA, alpha fades out
    f32   startSize    = 8.0f;
    f32   endSize      = 0.0f;
};
\end{lstlisting}

\begin{lstlisting}
ParticleEmitterComponent emitter;
emitter.maxParticles = 200;
emitter.emissionRate = 50.0f;
emitter.lifetime     = 0.8f;
emitter.velocityMin  = {-50.0f, -100.0f};
emitter.velocityMax  = {50.0f,  -200.0f};
emitter.startColor   = 0xFF8800FF;   // orange
emitter.endColor     = 0xFF880000;   // dim orange, fully transparent
world.add<ParticleEmitterComponent>(explosionEntity, emitter);
\end{lstlisting}

% ----------------------------------------------------------------------------
\section{3D Transform Components}

For 3D objects, use the components from \texttt{ecs/Components3D.hpp}.

\begin{lstlisting}
struct Position3D { Vec3 position; };
struct Rotation3D  { Vec4 quaternion = {0,0,0,1}; };  // identity
struct Scale3D     { Vec3 scale = {1,1,1}; };
\end{lstlisting}

These components store the local-space transform split into three separate
components instead of a single \texttt{Transform}. The scene system reads
them and computes the world matrix in \texttt{WorldTransform} (see
Chapter~\ref{chap:scene}).

\begin{lstlisting}
world.add<Position3D>(e, Position3D{{10.0f, 0.0f, 5.0f}});
world.add<Rotation3D>(e);           // identity rotation
world.add<Scale3D>(e, Scale3D{{2.0f, 2.0f, 2.0f}});
\end{lstlisting}
